% 20260626_Assingment_Probability_Statistics_A
\documentclass[dvipdfmx]{jsreport}
\usepackage{soupreamble}

\title{確率及び統計学特論A 第4回レポート課題}
\author{M1 6012 川村聡一郎}
\date{}


\begin{document}
\maketitle

\begin{tcolorbox}
\begin{enumerate}
	\item $a, b\in \R$に対して，$aX+ b$の特性関数は$\phi_{aX+ b}$とかけることを示せ．
	\item 確率変数$X, Y$は独立で，それぞれの特性関数を$\phi_X, \phi_Y$とする．このとき$X+ Y$の特性関数$\phi_{X+ Y}$は$\phi_X\phi_Y$とかけることを示せ．
	\item $X\sim \mathrm{Bin}(n, p)$のとき$X$の特性関数を求めよ\footnote{$\mathcal{P}(X= k)= \binom{n}{k}\cdot p^k\cdot (1- p)^{n- k},\ \ 0\le k\le n$}．
\end{enumerate}
\end{tcolorbox}

\begin{souanswer} %答案はここに書く
\begin{enumerate}
	\item dis:$aX+ b$とする．
\begin{align*}
	\phi_{aX+ b}(t)
	&= \mathbb{E}[\exp(it(aX+ b))]\\
	&= \mathbb{E}[\exp(itaX)\exp(itb)]\\
	&= \exp{(ibt)}\mathbb{E}[\exp(itaX)]\\
	&= \exp{(ibt)}\phi_X(at)\\
	&= \exp{(itb)}\phi_X(at)
\end{align*}

	\item dis:$X+ Y$とする．
\begin{align*}
	\phi_{X+ Y}(t)
	&= \mathbb{E}[\exp{(it(X+ Y))}]\\
	&= \mathbb{E}[\exp{(itX)}\cdot \exp{(itY)}]\\
	&= \mathbb{E}[\exp{(itX)}]\cdot\mathbb{E}[\exp{(itY)}]\\
	&= \phi_{X}\cdot \phi_{Y}
\end{align*}
	\item dis:$X\sim \mathrm{Bin}(n, p)$とする．
\begin{align*}
	\phi_{X\sim\mathrm{Bin}(n,p)}(t)
	&= \mathbb{E}\left[(e^{it})^k\cdot \binom{n}{k}\cdot p^k\cdot (1- p)^{n- k}\right]\\
	&= \mathbb{E}\left[(pe^{it})^k\binom{n}{k}\cdot (1- p)^{n- k}\right]\\
	&= \sum_{k=0}^{n} \binom{n}{k}\cdot (pe^{it})^k\cdot (1- p)^{n- k}\\
	&= (1- p+ pe^{it})^n
\end{align*}
\end{enumerate}
\end{souanswer}
\end{document}
